% =====================================================================
%  Thème 2 — Chasles et combinaisons — Niveau DIFFICILE — Corrigé
% =====================================================================
\documentclass[11pt,a4paper]{article}
\usepackage[utf8]{inputenc}
\usepackage[T1]{fontenc}
\usepackage{lmodern}
\usepackage[french]{babel}
\usepackage{amsmath,amssymb,mathtools}
\usepackage{xcolor}
\usepackage{enumitem}
\usepackage{fancyhdr}
\usepackage[a4paper,margin=2.1cm,headheight=26pt,headsep=10pt]{geometry}

\definecolor{primary}{RGB}{222,70,80}
\definecolor{primarydark}{RGB}{150,30,42}
\definecolor{excol}{RGB}{0,135,90}

\pagestyle{fancy}\fancyhf{}
\fancyhead[L]{\small\textcolor{primarydark}{\textbf{Fiche 2} — Calcul vectoriel}}
\fancyhead[R]{\small\textcolor{primarydark}{\textsc{Thème 2 — Difficile — Corrigé}}}
\fancyfoot[C]{\small\thepage}
\renewcommand{\headrulewidth}{0.8pt}

\newcommand{\vc}[1]{\overrightarrow{#1}}
\providecommand{\norm}[1]{\left\lVert #1\right\rVert}
\newcommand{\cor}[1]{\medskip\noindent{\color{primarydark}\bfseries\large Corrigé #1.}\par\smallskip}
\newcommand{\rep}{\textcolor{excol}{$\blacktriangleright$}\ }

\begin{document}
\begin{center}
{\LARGE\bfseries\color{primarydark}Chasles et combinaisons — Corrigé Difficile}\\[2pt]
{\color{primary}\rule{0.5\linewidth}{1.1pt}}
\end{center}
\vspace{2mm}

\cor{1}
\begin{enumerate}[label=\textbf{\alph*)},leftmargin=2em,itemsep=3pt]
  \item Sur la rangée du bas, $\vc{EF}$ et $\vc{FG}$ sont deux \og{}pas\fg{} d'une case vers la droite :
        même direction, même sens, même longueur, donc $\vc{FG}=\vc{EF}$. Par conséquent
        $\vc{FG}-\vc{EF}=\vec0$.
  \item $\vc{AB}+\vc{BF}=\vc{AF}$ (Chasles : la lettre $B$ s'efface).
  \item Il reste $\vec w=\vc{AF}+\vec0=\boxed{\vc{AF}}$.
\end{enumerate}
\rep On repère d'abord les termes qui \emph{se compensent} grâce à l'équipollence, puis on enchaîne le
reste par Chasles.

\cor{2}
\begin{enumerate}[label=\textbf{\alph*)},leftmargin=2em,itemsep=3pt]
  \item $\vc{MN}=\vc{MP}+\vc{PN}$. Comme $M$ est le milieu de $[PS]$, $\vc{MP}=-\dfrac12\vc{PS}$ ;
        comme $N$ est le milieu de $[PQ]$, $\vc{PN}=\dfrac12\vc{PQ}$. Donc
        \[\vc{MN}=\dfrac12\vc{PQ}-\dfrac12\vc{PS}.\]
  \item Dans le parallélogramme $PQRS$, les côtés $[PQ]$ et $[RS]$ sont opposés et parcourus en sens
        contraires, donc $\vc{PQ}=-\vc{RS}$.
  \item $2\,\vc{MN}=\vc{PQ}-\vc{PS}=-\vc{RS}-\vc{PS}$, soit
        \[\boxed{\,2\,\vc{MN}=-\vc{PS}-\vc{RS}\,}.\]
\end{enumerate}
\rep C'est exactement le mécanisme d'une question de concours : décomposer via un sommet commun ($P$),
puis remplacer $\vc{PQ}$ à l'aide des côtés opposés.

\cor{3}
\begin{enumerate}[label=\textbf{\alph*)},leftmargin=2em,itemsep=3pt]
  \item $ABCD$ parallélogramme $\iff \vc{AB}=\vc{DC}$. Or $\vc{AB}\,(5-1;3-2)=(4;1)$ et
        $\vc{DC}=(6-x_D\,;\,7-y_D)$. L'égalité donne $6-x_D=4$ et $7-y_D=1$, d'où
        $x_D=2$, $y_D=6$ : $\boxed{D\,(2;6)}$.
  \item $\vc{AD}\,(2-1;6-2)=(1;4)$ et $\vc{BC}\,(6-5;7-3)=(1;4)$ : on a bien $\vc{AD}=\vc{BC}$.
        Le quadrilatère est un parallélogramme.
  \item Le centre $I$ est le milieu commun des deux diagonales. Milieu de $[AC]$ :
        $I\Bigl(\dfrac{1+6}{2};\dfrac{2+7}{2}\Bigr)=\Bigl(\dfrac72;\dfrac92\Bigr)=(3{,}5\,;\,4{,}5)$.
        \emph{(Vérification par $[BD]$ : $\bigl(\frac{5+2}{2};\frac{3+6}{2}\bigr)=(3{,}5;4{,}5)$.)}
\end{enumerate}
\rep Traduire une propriété géométrique (parallélogramme) par une \emph{égalité vectorielle}, puis
passer aux coordonnées : c'est le pont entre le Thème 1 et le Thème 2.

\end{document}
